\section{Abstraction}\label{sec:lit_rev:abstraction}
Abstraction is a fundamental concept in computer science and software engineering,providing the basis for tooling and development in both fields.
Abstraction is equally fundamental as a mode of thinking in other domains,describing an approach that is useful for separating concerns and reducing complexity~\cite{Wing2006ComputationalThinking}.
Examples of this cross-domain application appear in fields such as mathematics <TODO: example and citation>, physics <TODO: example and citation> and linguistics <TODO: example and citation>.

Good abstraction provides practitioners of any field with the ability to remove complexity without losing expressivity in the solutions they can explore and produce~\cite{parnas1972criteria,aho2022abstractions}.
A well-designed abstraction reduces complexity for its user through well-defined interfaces, without constraining the space of solutions they can express or construct~\cite{parnas1972criteria}.

As this section demonstrates, quantum computing has not yet realised the abstraction quality that characterises mature classical programming systems: existing approaches improve ergonomics but largely operate at the functional level, leaving correctness, portability and conceptual expressivity underserved.
Establishing a principled, concrete definition of abstraction and a framework for differentiating abstraction levels is therefore a prerequisite for any rigorous evaluation of quantum programming systems.
This section develops that baseline, drawing on classical computing literature, and applies it to characterise the current state of abstraction in quantum computing.

The resulting framework is carried forward throughout the remainder of this literature review and underpins the critical analysis of quantum programming language abstractions presented in Chapter~\ref{chapt:pub_language_review}.

\subsection{Defining Abstraction}\label{subsec:lit_rev:abstraction:defining_abstraction}
While abstraction is useful both within and beyond computer science, this subsection focuses on its definition specifically in computer science and software engineering.
The aim is to present the ways in which abstraction is defined and discussed in the literature and to identify the roles and uses of different definitions.
These definitions ground the discussion throughout the remainder of this section.
The focus is predominantly on literature published since 2020, though older works establishing standard interpretations of foundational concepts are drawn upon where appropriate.

Abstraction as used in this subsection may refer to either a layered system with reduced or increased complexity between layers or to a specific layer within such a system.
This discussion distinguishes the two using the terms \textit{System (of abstraction)} and \textit{Abstraction (layer)} respectively.
Regardless of which form is under discussion, any system of abstraction and the layers within it share a common set of defining elements:
\begin{itemize}
    \item Decomposition of complexity into generalised tractable domains: assumptions delimit the scope of the system, allowing complexity to be managed without exhaustive enumeration~\cite{Wing2006ComputationalThinking}
    \item Presence of data models and operations over them: the system defines what it acts on and how~\cite{aho2022abstractions}
    \item A mechanism for transitioning between more and less general layers: complexity can be recovered or removed as required~\cite{aho2022abstractions}
\end{itemize}
These three criteria are present in any system of abstraction.

The presence of greater generality or more assumptions is commonly termed \textit{higher level}; lesser generality and greater specificity, \textit{lower level}.
In computer science, compilation refers to transitions between levels of generality within a system of abstractions~\cite{aho2022abstractions}.

Higher- to lower-level relationships also arise between distinct systems of abstraction, where a higher-level system is grounded in or augmented by a lower-level description of the same domain.
An example of this relationship between abstraction systems is the three-layered model of parallel computing~\cite{DiMatteo2024AbstractionHierarchy}, comprising the \textbf{Programming}, \textbf{Execution} and \textbf{Hardware} layers.
Each layer makes assumptions about the availability of operations in the layer below, thereby increasing its own generality; Figure~\ref{fig:lit_rev:abstraction:abstraction_relationships} describes these dependency relationships.

Additional specificity or generality can be introduced on either side of this layered structure: at the more general end, all layers may collapse into a single bridging model~\cite{Valiant1990BSP} used as its own layer in higher-level systems.
At the more specific end, the programming layer may be expanded to include high-level languages such as Python or domain specifications for machine learning.

\begin{figure}[h!]
    \centering
    \includegraphics[width=0.75\linewidth]{../figures/fig:lit_rev:abstraction:abstraction_relationships}
    \caption{A representation of the common programming, execution, hardware abstraction system with concrete examples for a C program.}
    \label{fig:lit_rev:abstraction:abstraction_relationships}
\end{figure}

The scope of any system of abstraction is therefore contingent on the discussion it is designed to support; the definitions established here are applied with that flexibility in mind.

It is necessary to be able to compare both systems of abstraction and abstractions to determine their applicability to a given domain and value in a situation.
Comparative relationships between abstractions are well documented; however, formalisms for discussing such relationships between systems of abstraction are absent from the literature.
The taxonomy of abstractions~\cite{aho2022abstractions} provides a spectrum from Fundamental → Abstract → Declarative → Computational, offering a framework for describing relationships between abstractions in computational processes.
Through this taxonomy it is possible to relate abstractions at higher or lower levels to each other, forming systems of abstractions.
As an example, Table~\ref{tab:lit_rev:abstraction:abstraction_taxonomy} illustrates the taxonomic relationships of dictionaries and maps~\cite{aho2022abstractions}.

\begin{table}[htbp]
    \centering
    \caption{Taxonomy of abstractions illustrating the progression from fundamental concepts to computational implementations for both classical and quantum domains, adapted from Aho and Ullman~\cite{aho2022abstractions}.}
    \label{tab:lit_rev:abstraction:abstraction_taxonomy}
    \begin{tabular}{@{}lll@{}}
        \toprule
        \textbf{Taxonomy Level} & \textbf{Classical Example}          \midrule
        Fundamental & Dictionary (insert, delete, lookup) \\
        Abstract & LinkedList, Map \\
        Declarative & SQL Table/Query \\
        Computational & Java/C implementation & QASM \\ \bottomrule
    \end{tabular}
\end{table}

Comparative frameworks such as the taxonomy are useful for comparing abstractions but provide no real measure of comparison between systems.
Two complementary definitions are introduced here that will allow us to discuss comparisons between more and less abstract systems.
A system could simultaneously belong to both definitions based on its position within the spectrum of a specific domain, as we saw with parallel computing.
\begin{itemize}
    \item \textit{Conceptual abstraction}: An abstraction characterised by a small number of broad, inherent concepts that define its operational scope.
    \item \textit{Functional abstraction}: A hierarchy of abstractions forming layers with distinct implementation roles.
\end{itemize}

These definitions provide conceptual tools for comparing abstractions, but do not suffice for critically examining the state of abstraction in quantum computing addressed in Subsection~\ref{subsec:lit_rev:abstraction:in_quantum_computing}.
The taxonomy describes abstraction across computation broadly and within systems; it does not provide the vocabulary to locate languages within a compilation stack specifically.
The system comparison framework likewise allows us to reason about the applicability of different systems within a domain but does not define which system configurations serve compilation or computation.
More precise vocabulary is therefore required that can:
\begin{itemize}
    \item Define systems of abstraction for programming language compilation.
    \item Categorise and place languages within systems of abstraction.
\end{itemize}

These tools will allow us to understand the roles of different quantum programming languages discussed in Section~\ref{sec:lit_rev:quantum_programming_languages} and develop the assessment frameworks in Chapter~\ref{chapt:pub_language_review}.
To arrive at our required definitions, we first need to explore the ways in which abstraction is used in programming.

The journey to define abstraction within programming and computer science has been ongoing since the 20th century.
The development of higher level programming languages and increasingly complex systems required definitions to delineate how and when abstraction should and could be used.

Abstraction in programming defines more than just how languages within a compilation stack relate to each other, it also defines how systems are designed and modularised.
The language used to describe both is rooted in the same general principles:
\begin{itemize}
    \item Abstraction hides information and defines boundaries of responsibilities~\cite{parnas1972criteria}
    \item Abstraction preserves semantic intention between layers~\cite{yuan2025foundationalabstractions}
    \item Abstraction improves a user's ability to reason with the layer of a system~\cite{DiMatteo2024AbstractionHierarchy}
\end{itemize}

In programming languages abstraction defines the responsibility of a language to both the user and the machine.
The system of abstraction selected for a given domain must responsibly hide information at each layer without inhibiting the user's ability to express their intent.
Simultaneously it must ensure that intention is preserved so that the machine can execute it faithfully.

But abstraction also defines the limitations of control a user of a language has on the system they are manipulating~\cite{yuan2025foundationalabstractions}.
Abstraction in programming removes responsibility over details or actions that can be assumed or hidden, which also removes the ability to control them.
A canonical example in classical programming is the presence or absence of garbage collection in a given language.

Now can it be said that Java is a higher abstraction of C given one implements garbage collection and the other doesn't?
Most definitions of abstraction in programming would acknowledge Java is higher-level but wouldn't necessarily place it at a higher \textit{abstraction layer} within a programming system.
Why is that?

Abstraction in programming is defined by more than just information hiding or reduction in control.
In programming intention influences abstraction, the intention of a language in concert with the actions or control it permits is what places it at a certain layer in a system.
Languages like C or Java sit at the same layer in parallel systems of abstraction because they inhabit the same intention, for example, application development.
It's for this reason we don't look at C as being an assembly language despite its ability to have low-level control over memory; that's not its intended purpose.

The preceding discussion motivates the following definitions, which are used throughout this review.

\textit{Language abstractions} are user-facing representations that hide implementation details from an end user.
Implementation details may include specifics for interfacing with a machine, physical system properties or operations whose behaviour can be inferred from the domain context.
A language abstraction promotes tractability within a domain by improving a user's ability to reason about the behaviour of programs in that domain.
A well-formed language abstraction preserves expressivity within the domain, even where it constrains execution-level control or reduces access to low-level implementation choices.
For example, a user need not specify the underlying character array representation of a string yet retain full expressive capacity for string manipulation.

\textit{Compiler chains} define computational systems of abstraction.
Compiler chains link high-level to low-level language abstractions by translating from more to less general representations.
A well-formed compiler chain guarantees no loss of intention across layers; higher layers are expected to reduce the cognitive burden on the programmer.
Compiler chains can be compared faithfully by expressing the same system intention in their layers.

The following section examines how these chains have been realised in classical computing systems, establishing a comparative baseline for evaluating quantum programming abstractions later in this review.

\subsection{Levels and elements of abstraction}\label{subsec:lit_rev:abstraction:levels_elements}
Abstraction is defined in multiple ways as shown in Section~\ref{subsec:lit_rev:abstraction:defining_abstraction}.
Abstraction can be defined relatively between layers in a system using the taxonomy of abstraction~\cite{aho2022abstractions} or relatively between systems.
In computation abstraction is defined both as languages that inhabit layers and compiler chains that describe systems of layers.
This subsection focuses specifically on those systems.

Examples will be presented of different ways of organising abstraction within classical computing domains and a framework for assessing the effectiveness of those abstractions provided.
Figure~\ref{fig:lit_rev:abstraction:classical_compiler_chains} provides a graphical representation illustrating three chains of compilation within classical systems.
The first is one most software engineers and programmers would be familiar with; C may be replaced with many high-level languages and the structure would still hold.
The C chain describes a common separation of layers in programming abstraction with programming, execution and hardware layers.
The other two examples build on this with the SQL-led chain showing an extended compiler chain with a domain specification language (DSL) included.
The Python chain provides a discussion point.

Based on Figure~\ref{fig:lit_rev:abstraction:classical_compiler_chains} Python is serving as a machine learning (ML) DSL as it has come to be used in recent years~\cite{agrawal2019tensorflow}.
However, Python can also be a scripting or programming layer; this is an example of how the compilation chains are flexible based on the intention of a layer.
What matters here is the audience of the abstraction and how the choice of abstraction benefits or hinders their ability to interact with it~\cite{DiMatteo2024AbstractionHierarchy}.

\begin{figure}[h!]
    \centering
    \includegraphics[width=0.75\linewidth]{../figures/fig:lit_rev:abstraction:classical_compiler_chains}
    \caption{An example of three different ways to interpret compiler chains showing a traditional application development chain, SQL database domain chain and Python's use in machine learning.}
    \label{fig:lit_rev:abstraction:classical_compiler_chains}
\end{figure}

When Python acts as an ML DSL it provides the purpose of hiding common programming paradigms and simplifying the way ML engineers and data scientists interact with machines.
SQL serves a similar purpose allowing non-programmers to interact with computational systems in a way that benefits their domain and level of understanding.
What counts as the highest abstraction level is relative to the domain and intended audience; in both cases presented the audience loses control over memory management or algorithmic decisions.
Despite these losses, the information they hide allows the user to experience a more tractable version of computation in line with their understanding.

The layers of a compiler chain can be described as in Table~\ref{tab:lit_rev:abstraction:compiler_layers}.

\begin{table}[htbp]
    \centering
    \caption{General layers of a compilation chain including domain-specific simplification for high-level user interaction, demonstrates the cognitive simplification and distance from machine representations influence layering.}
    \label{tab:lit_rev:abstraction:compiler_layers}
    \begin{tabular}{@{}llp{8cm}@{}}
        \toprule
        \textbf{Layer} & \textbf{Level} & \textbf{Primary Purpose}                                                             \\ \midrule
        Domain         & Highest        & Hiding programming paradigms; providing domain-specific syntax (e.g., SQL, ML DSLs). \\
        Source         & High           & Human-readable logic; memory and resource management abstractions.                   \\
        Intermediate   & Middle         & Machine-independent optimization and canonical representation (IR).                  \\
        Assembly       & Low            & Mapping logic to specific hardware instructions and registers.                       \\
        Binary         & Lowest         & Execution on physical hardware.                                                      \\ \bottomrule
    \end{tabular}
\end{table}

While these layers appear well-structured, whether these compilation chains are effective needs to be qualified.
Good abstraction must match the intention of the definition outlined in Subsection~\ref{subsec:lit_rev:abstraction:defining_abstraction}.
To recap, that means:
\begin{itemize}
    \item Abstraction hides information and defines boundaries of responsibilities~\cite{parnas1972criteria}
    \item Abstraction preserves semantic intention between layers~\cite{yuan2025foundationalabstractions}
    \item Abstraction improves a user's ability to reason with the layer of a system~\cite{DiMatteo2024AbstractionHierarchy}
\end{itemize}

Effective abstractions guarantee that definition and can be assessed by the following criteria:
\begin{itemize}
    \item Correctness of execution can be verified through conditions and invariants~\cite{yuan2025foundationalabstractions, LiskovGuttag1986}
    \item Data types define what they do, not how they do it~\cite{LiskovGuttag1986}
    \item Procedures have formal, guaranteed definitions~\cite{LiskovGuttag1986}
    \item Produced solutions reduce in complexity in higher layers~\cite{aho2022abstractions,DiMatteo2024AbstractionHierarchy}
    \item Layers can be internally changed without \textit{requiring} changes in other layers~\cite{parnas1972criteria}
\end{itemize}

Table~\ref{tab:lit_rev:abstraction:assessing_sql_compiler_chain} applies these criteria to the SQL compiler chain in Figure~\ref{fig:lit_rev:abstraction:classical_compiler_chains}.
\begin{table}[htbp]
    \centering
    \caption{Assessment of the SQL compiler chain based on effective abstraction criteria.}
    \label{tab:lit_rev:abstraction:assessing_sql_compiler_chain}
    \begin{tabular}{@{}lp{10cm}@{}}
        \toprule
        \textbf{Criterion}    & \textbf{SQL Implementation Evidence}                                                      \\ \midrule
        Information Hiding    & Users interact with relations/tables; physical storage and indexing are hidden.           \\
        Semantic Preservation & Declarative intent (WHAT) is preserved through optimization to relational algebra (HOW).  \\
        Reasoning Improvement & Non-programmers can query complex datasets without understanding B-trees or buffer pools. \\
        Formal Definition     & Relational algebra provides a formal, guaranteed basis for procedures.                    \\
        Reduced Complexity    & High-level syntax reduces the cognitive load compared to manual data retrieval loops.     \\ \bottomrule
    \end{tabular}
\end{table}

Understanding what an abstraction is and how to assess it provides an insight into why abstraction exists and its role in simplifying complex systems.
What these concepts do not provide is commentary on why it is important for the users of those systems or why it matters to the design of quantum computing languages.

The criteria established here, information hiding, semantic preservation, formal definition, and complexity reduction will serve as the evaluative baseline when comparing quantum programming systems in the sections that follow.

In Subsection~\ref{subsec:lit_rev:abstraction:why_matters} a review concerning the importance of traditional and non-traditional abstraction systems is presented.
The review provides qualitative results on why it is important to continue to examine abstraction in emerging systems and the benefits it provides their users.

%% ------------------------ To be written --------------------------------------------------------------------------- %%

\subsection{Why abstraction matters}\label{subsec:lit_rev:abstraction:why_matters}
- Computational thinking and how it is described by languages
- Incorporate natural language to SQL article
- Incorporate article on phone number empirical comparison
- Look at adoption statistics for LLM based systems
- Accessibility of systems with higher levels of abstraction
- aho2022 will be important her as well for historical relevance and importance of abstraction
- What makes an abstraction worthwhile may be worth a cursory mention here to justify why it matters e.g abstraction when done well looks like A, this is useful because of B
- Discuss tractability from Wing 2006 and link it back to findings from the quantum survey

\subsection{Abstraction in quantum computing}\label{subsec:lit_rev:abstraction:in_quantum_computing}
- Discuss DiMatteo here and begin to disect why that abstraction framework isn't enough
- Need to nail down here the failures of abstraction in quantum computing, di matteo can help with this
- Is conceptual abstraction missing or too low level? Maybe comment on the prevalence of functional abstractions

\subsection{Transition (TBC)}\label{subsec:lit_rev:abstraction:transition}
- Pull together the points and findings of the other subsections here to justify that abstraction is the core missing piece of quantum computing
- Make a point that we are missing a level of abstraction by comparing the presence of python and SQL in classical languages
- Provide bridging here that we must focus through that lens for the rest of the review and how it will motivate the thesis.
